% SPDX-License-Identifier: CC-BY-NC-ND-4.0
% Copyright 2026 Ingolf Lohmann.
\documentclass[10pt,a4paper]{article}
\usepackage[a4paper,top=16mm,bottom=17mm,left=18mm,right=18mm,headheight=13pt]{geometry}
\usepackage{fontspec}
\setmainfont{P052-Roman.otf}[
  Path=/usr/share/fonts/opentype/urw-base35/,
  BoldFont=P052-Bold.otf,
  ItalicFont=P052-Italic.otf,
  BoldItalicFont=P052-BoldItalic.otf]
\setsansfont{DejaVu Sans}
\usepackage[provide=*,ngerman]{babel}
\usepackage{microtype,xcolor,graphicx,booktabs,tabularx,enumitem,tikz}
\usetikzlibrary{arrows.meta,positioning}
\usepackage[most]{tcolorbox}
\usepackage{fancyhdr,hyperref,multicol}
\definecolor{Ink}{HTML}{152B38}
\definecolor{Ocean}{HTML}{147D92}
\definecolor{Mint}{HTML}{27A58B}
\definecolor{Sun}{HTML}{E0A32B}
\definecolor{Mist}{HTML}{ECF6F7}
\hypersetup{colorlinks=true,urlcolor=Ocean,linkcolor=Ocean,
 pdftitle={Ein Spiegel für Ursachen: QIK-VRT},pdfauthor={Ingolf Lohmann}}
\pagestyle{fancy}\fancyhf{}
\fancyhead[L]{\sffamily\small Ein Spiegel für Ursachen}
\fancyhead[R]{\sffamily\small QIK-VRT · 2026}
\fancyfoot[C]{\sffamily\small\thepage}
\renewcommand{\headrulewidth}{0.3pt}
\setlist[itemize]{leftmargin=1.3em,itemsep=2pt}
\setlength{\emergencystretch}{2em}
\newtcolorbox{keybox}{enhanced,breakable,colback=Mist,colframe=Mint,boxrule=.8pt,
 arc=2mm,left=4mm,right=4mm,top=3mm,bottom=3mm}
\newtcolorbox{limitbox}{enhanced,breakable,colback=yellow!6,colframe=Sun,boxrule=.8pt,
 arc=2mm,left=4mm,right=4mm,top=3mm,bottom=3mm}
\newcommand{\head}[1]{\vspace{4mm}{\sffamily\Large\bfseries\color{Ink}#1}\par\vspace{1.5mm}}
\begin{document}
\thispagestyle{empty}
\begin{tikzpicture}[remember picture,overlay]
\fill[Ink] (current page.north west) rectangle ([yshift=-82mm]current page.north east);
\fill[Ocean] ([xshift=132mm]current page.north west) -- (current page.north east) --
 ([yshift=-82mm]current page.north east) -- ([xshift=92mm,yshift=-82mm]current page.north west) -- cycle;
\end{tikzpicture}
\vspace*{9mm}\color{white}
{\sffamily\small\bfseries\color{Sun} WISSENSCHAFT · INFORMATIK · PHYSIK\par}
\vspace{8mm}
{\sffamily\fontsize{29}{34}\selectfont\bfseries Ein Spiegel\\für Ursachen\par}
\vspace{5mm}
{\Large Wie QIK-VRT aus Daten einen\\überprüfbaren Erkenntnisweg macht\par}
\vspace{8mm}{\sffamily Ingolf Lohmann\par}
\vspace{19mm}\color{black}

\begin{keybox}
Computer finden Informationen rasend schnell. Die schwierigere Frage lautet:
Warum ist eine konkrete Aussage belastbar? QIK-VRT bindet Quelle, Messung,
Transformation, Hypothese, Beweis, Widerspruch und Wirkung zu einem
überprüfbaren Graphen. 21 Lean-Sätze sichern den endlichen Strukturkern ab.
Ein universeller Wahrheitsautomat entsteht ausdrücklich nicht.
\end{keybox}

\begin{multicols}{2}
Ein Suchtreffer kann eine Behauptung wiederholen. Ein Sprachmodell kann sie
flüssig erklären. Ein Sensor liefert Zahlen, ein Beweisassistent akzeptiert
Formeln. Aber keiner dieser Vorgänge enthält automatisch den ganzen
wissenschaftlichen Weg. Oft fehlen Kalibrierung, Prämissen, Unsicherheit,
Gegenbelege oder die Rückmeldung, welche Wirkung tatsächlich eintrat.

QIK-VRT lässt sich als \textbf{Kausalitätsspiegel} lesen. Gespeichert wird nicht
nur das Ergebnis, sondern der nachvollziehbare Weg von Beobachtung über Modell
und Entscheidung bis zur beabsichtigten und beobachteten Wirkung. Der Spiegel
ist nicht die Wirklichkeit. Er zeigt, welche Brücke belegt, angenommen oder
noch offen ist.

\head{Sechs Erkenntnisarten}

Statt eines einzigen Wahrheitsbits unterscheidet das Protokoll: formal
bewiesen, empirisch gestützt, quellengebunden, normativ, interpretativ und
offen. Ein Kernelbeweis gilt im genannten Modell. Eine Messung gilt in ihrem
Mess- und Unsicherheitsrahmen. Eine Quelle belegt, wer etwas sagte, nicht dass
es stimmt. Eine offene Aussage darf nicht durch sprachliche Sicherheit zum
Fakt werden.

Diese Typisierung verhindert Kategoriefehler. Lean prüft keine Sensoren. Ein
Messwert beweist keinen Allgemeinsatz. Eine DOI garantiert eine fixierte
Ablage, nicht wissenschaftliche Wahrheit. Und eine technisch berechnete
Empfehlung ist keine Erlaubnis, eine Maschine zu bewegen.

\head{Was Lean geprüft hat}

Lean 4.19.0 kompilierte 21 Sätze ohne ausgelassene Beweise und ohne
projektspezifische Axiome. Bewiesen sind unter anderem:

\begin{itemize}
\item Additive Erweiterung erhält vorhandene Objekte.
\item Mesh-Vereinigung ist kommutativ, assoziativ und idempotent.
\item Replikate konvergieren nach denselben Updates unter gleicher Policy.
\item Vorhandene Begründungspfade und Konflikte bleiben sichtbar.
\item Exakte Neuheit ist nur relativ zum untersuchten Korpus.
\item Derselbe Trace kann verschiedene physische Ursachenzuschreibungen tragen.
\item Digital-Twin-Aktorwirkung braucht qualifizierte Beobachtung und Effect-Ack.
\item Proposal-only Analyse autorisiert keine Außenwirkung.
\end{itemize}

Das Gegenbeispiel ist ebenso wichtig: Ein leerer Korpus beantwortet keine
Frage. Die Struktur garantiert also nicht, jede Frage beantworten zu können.

\head{Von 94 Grad zur Ursache}

Ein Sensor meldet 94 Grad, ein Modell prognostiziert einen Ausfall, ein Mensch
bestätigt die Abschaltung, der Aktor reagiert, die Temperatur sinkt. Eine
Logdatei zeigt Zeitstempel. Der Kausalitätsspiegel fragt zusätzlich: Welcher
Sensor? Welche Kalibrierung und Unsicherheit? Welche Modellversion? Welche
Prämissen? Empfehlung oder ausgeführte Aktion? Welche Rückmessung? Welche
alternative Ursache?

So wird aus einer Geschichte ein prüfbarer Graph. Physische Ursache folgt
dennoch nicht aus Chronologie. Je nach Frage braucht man Intervention,
Kontrollgruppe, Identifikationsannahme oder validiertes Domänenmodell. QIK-VRT
macht den fehlenden Brückenschritt sichtbar.

\head{Analoge Welt und Digital Twin}

Ein analoges Signal wird durch Sensor, Übertragungsfunktion, Abtastung,
Quantisierung und Zeitstempel zu digitalen Daten. QIK-VRT kann Messwert,
Geräteidentität, Kalibrierparameter, Unsicherheitsmodell und Auswertecode
zusammenbinden. Korrekturen werden neue, rückverweisende Knoten.

Beim Digital Twin werden physisches Objekt und laufendes Modell gekoppelt.
Gerade Übergänge sind fehleranfällig: falsche Einheit, veraltete Modellfassung,
fehlende Kalibrierung, Datenlücke oder eine Empfehlung, die wie eine
ausgeführte Handlung behandelt wird. Explizite Kanten unterstützen
Root-Cause-Analyse, Wartung, Validierung und Audit. Sie ersetzen keine
anlagenspezifische Sicherheitsprüfung.

\head{Quantenebene ohne Magie}

Ein Quantenexperiment besitzt eine lange klassische Kette: Schaltkreis,
Compiler, Backend-Kalibrierung, Shots, Statistik und Interpretation. QIK-VRT
kann jeden Schritt content-addressiert binden. Auch Modelle ohne fest
vorgegebene globale Kausalordnung lassen sich dokumentieren, ohne daraus
voreilig einen Nachrichtenkanal in die Vergangenheit zu machen.

Der Wert der Quantenkausalebene ist die genaue Trennung von mathematischem
Modell, Versuchsaufbau, Statistik und physikalischer Deutung. Nicht bewiesen
sind VRT-Emergenz, Quanten-zu-Klassik-Limes, physische Retrokausalität oder ein
realer QPU-End-to-End-Kanal.

\head{Ein wachsender Erkenntnisgraph}

Inhaltsadressierte Repositories können Objektmengen stabil vereinigen. Ein
Widerspruch wird nicht überschrieben, sondern samt Quelle und Methode
erhalten. Dadurch kann Suche nicht nur nach ähnlichen Wörtern fragen, sondern
nach Abhängigkeiten, Widerlegungen, Kernelversionen und vollständigen
Evidenzpfaden.

Eine Antwort kann als kleiner Begründungsgraph ausgegeben werden. Das kann
Menschen und künstlich-kognitiven Systemen helfen, weil nicht nur das Ergebnis,
sondern sein Weg überprüfbar ist. Ob dadurch alle Nutzenden kognitiv besser
werden, bleibt eine empirische Forschungsfrage mit Messgrößen,
Vergleichsgruppen und Replikation.

\head{Wert für die Fachwelt}

Mathematiker erhalten einen versionierten Beweisgraphen: Definition,
Prämissen, exakte Satzfassung und Kernel-Receipt bleiben verbunden. Physiker
können Rohdaten, Kalibrierung, Code, Unsicherheit und konkurrierende Modelle
trennen. Informatik und Nachrichtentechnik gewinnen ein deterministisches
Evidenz- und Merge-Protokoll. Mess- und Regeltechnik erhalten eine
auditierbare Schleife von Sensor bis Wirkung.

\head{Verantwortung im Protokoll}

Klassifikation, Review, Repository-Merge, Zenodo-Ablage, IETF-Einreichung und
reale Aktorwirkung sind getrennte Effekte. Diese Trennung unterstützt auch
AI-Audits: Welche Daten und Modellfassung erzeugten eine Ausgabe? Welche
Unsicherheit war bekannt? Wer gab frei? Welche Wirkung wurde beobachtet?
QIK-VRT schafft Evidenzbereitschaft, aber keine automatische Zertifizierung
oder Rechtsberatung.

\end{multicols}

\begin{limitbox}
\textbf{Die präzise Aussage:} Im angegebenen endlichen Modell können
Erkenntnisobjekte statusgebunden wachsen, Konflikte erhalten bleiben,
Mesh-Replikate deterministisch vereinigen und Außenwirkungen getrennt
autorisiert werden. Universelle Wahrheit, globale Neuheit, Antworten auf jede
Frage und eine physische Quanten- oder Retrokausalitätsbrücke werden nicht
beansprucht.
\end{limitbox}

\vfill
{\sffamily\small\color{Ink}
\textbf{Quelle und Reproduzierbarkeit:} Goldkelch/qik-vrt.\\
Publikationspaket \texttt{2026-08-01-scientific-fact-growth-mesh}.\\
Additiv zum Zenodo-Addendum
\href{https://doi.org/10.5281/zenodo.21731492}{doi:10.5281/zenodo.21731492}.
Autor: Ingolf Lohmann. Geltungsgrenzen und Kernel-Receipt sind Bestandteil des
Pakets.}
\end{document}
